\chapter{DFS recursivo: el \emph{baseline}}
\label{cap:dfs-baseline}

\begin{intuicion}
Antes de hablar de \emph{worst-case-optimal joins} y semi-join
reducers, conviene tener firme el algoritmo más básico: DFS
(depth-first search) recursivo. Es el algoritmo \emph{baseline} del
matcher; cuando ningún operador especializado dispatcha, el motor
cae a un DFS bien podado. Este capítulo lo presenta primero en su
forma genérica (sobre grafos sin etiquetas, sin tiempo) y después
adaptado al subgraph matching temporal etiquetado de GraphHunter.
Si DFS te resulta familiar, podés saltar directo a la
sección~\ref{sec:dfs-temporal}; si no, esto es el refresher.
\end{intuicion}

\section{DFS clásico sobre un grafo}

\begin{definition}[DFS]
Dado un grafo dirigido \(G = (\V, \E)\) y un vértice de partida
\(s \in \V\), DFS visita todos los vértices alcanzables desde \(s\)
explorando \emph{en profundidad}: avanza por una arista, recursa al
nuevo vértice, y sólo retrocede cuando todas las salidas del vértice
actual ya fueron exploradas.
\end{definition}

\begin{algorithm}[H]
\caption{DFS recursivo, versión \emph{visit-all}}
\KwIn{grafo \(G = (\V, \E)\), vértice de partida \(s\)}
\KwOut{conjunto de vértices visitados}
\(\mathrm{visited} \leftarrow \emptyset\)\;
\(\mathrm{visit}(s)\)\;
\BlankLine
\Comment{Procedimiento \texttt{visit(v)}}
\If{\(v \in \mathrm{visited}\)}{\Return}
\(\mathrm{visited} \leftarrow \mathrm{visited} \cup \{v\}\)\;
\For{\(w \in N^{+}(v)\)}{
    \(\mathrm{visit}(w)\)\;
}
\end{algorithm}

\begin{ejemplo}
DFS sobre un grafo de 5 vértices arrancando desde \(A\).

\begin{figure}[H]
\centering
\begin{tikzpicture}[
    every node/.style={draw=accent, thick, circle, minimum size=8mm,
                       fill=accentlight, font=\small},
    edge/.style={->, >=Latex, accent, thick}
]
\node (A) at (0,2) {A};
\node (B) at (-1.5,0.5) {B};
\node (C) at (1.5,0.5) {C};
\node (D) at (-1.5,-1) {D};
\node (E) at (1.5,-1) {E};

\draw[edge] (A) -- (B);
\draw[edge] (A) -- (C);
\draw[edge] (B) -- (D);
\draw[edge] (B) -- (E);
\draw[edge] (C) -- (E);
\end{tikzpicture}
\caption{Grafo de ejemplo. DFS desde \(A\): \(A \to B \to D\)
(retroceso) \(\to E\) (retroceso a \(B\) y a \(A\)) \(\to C\)
(\(E\) ya visited, sin más). Orden de visita: \(A, B, D, E, C\).}
\label{fig:dfs-ejemplo}
\end{figure}
\end{ejemplo}

\subsection{Complejidad}

DFS visita cada vértice y arista una vez:
\begin{itemize}
    \item Tiempo: \(O(|\V| + |\E|)\). Lineal en el tamaño del grafo
    alcanzable.
    \item Espacio: \(O(|\V|)\) por la \texttt{visited} más
    \(O(\text{depth})\) por la pila de recursión.
\end{itemize}

\section{De \emph{visit-all} a \emph{enumerate-paths}}

DFS \emph{visit-all} visita cada vértice una vez. Para subgraph
matching necesitamos algo distinto: enumerar \emph{caminos}
(secuencias de vértices) que matchean un patrón. La diferencia es
crítica:

\begin{itemize}
    \item En \emph{visit-all}, el vértice se marca visited y no se
    explora de nuevo aunque haya otra forma de llegar a él.
    \item En \emph{enumerate-paths}, queremos \emph{todas} las formas
    de llegar a un vértice (cada una es un path distinto), pero
    \textbf{sin repetir vértices dentro de un mismo path}
    (k-simplicidad, capítulo~\ref{cap:grafos-temporales}).
\end{itemize}

\begin{algorithm}[H]
\caption{DFS \emph{enumerate-paths}, longitud exacta \(k\)}
\KwIn{grafo \(G\), vértice \(s\), longitud \(k\)}
\KwOut{secuencia de paths de longitud \(k\) desde \(s\)}
\(\mathrm{path} \leftarrow [s]\)\;
\(\mathrm{recurse}(s, k)\)\;
\BlankLine
\Comment{\texttt{recurse(v, remaining)}}
\If{\(\mathrm{remaining} = 0\)}{
    \texttt{emit}(\(\mathrm{path}\))\;
    \Return\;
}
\For{\(w \in N^{+}(v)\)}{
    \If{\(w \notin \mathrm{path}\)}{
        \(\mathrm{path}.\texttt{push}(w)\)\;
        \(\mathrm{recurse}(w, \mathrm{remaining} - 1)\)\;
        \(\mathrm{path}.\texttt{pop}()\)\;
    }
}
\end{algorithm}

\subsection{Backtracking}

Notar el patrón \texttt{push} / \texttt{recurse} / \texttt{pop}: es
el \emph{backtracking} clásico. La pila \texttt{path} se modifica
durante la recursión y se restaura al volver, manteniendo el
estado consistente para la próxima rama.

\begin{nota}
El chequeo \(w \notin \mathrm{path}\) implementa la
\emph{1-simplicidad} (path simple). Para \(k\)-simplicidad general,
contamos cuántas veces aparece \(w\) en \texttt{path} y aceptamos si
es \(< k\). En GraphHunter el chequeo está en
\filepath{core/graph-engine/src/dfs.rs}: cuenta apariciones via
\emph{linear scan} sobre el binding (siempre tiene
\(\leq k+1\) elementos, costo despreciable).
\end{nota}

\section{Subgraph matching: agregar etiquetas}

Para subgraph matching del capítulo~\ref{cap:subgraph-matching}, DFS
necesita ajustes:

\begin{enumerate}
    \item Sólo visitar vértices del tipo correcto en cada paso.
    \item Sólo seguir aristas del tipo correcto en cada paso.
    \item Acumular el path \emph{etiquetado} para emit.
\end{enumerate}

\begin{algorithm}[H]
\caption{DFS \emph{subgraph match}, hipótesis chain
\((L_0^V, L_0^E, L_1^V, L_1^E, \dots, L_k^V)\)}
\KwIn{grafo \(G\), hipótesis \(H\)}
\KwOut{Vec<HuntResult>}
\For{\(s \in \V\) con \(\mathrm{type}(s) = L_0^V\)}{
    \(\mathrm{path} \leftarrow [s]\)\;
    \(\mathrm{recurse}(s, 0)\)\;
}
\BlankLine
\Comment{\texttt{recurse(v, step)}}
\If{\(\mathrm{step} = k\)}{
    \texttt{emit}(\(\mathrm{path}\))\;
    \Return\;
}
\For{\(w \in N^{+}(v)\)}{
    \(e \leftarrow (v, w)\)\;
    \If{\(\mathrm{type}(e) = L_{\mathrm{step}}^E\) \textbf{and} \(\mathrm{type}(w) = L_{\mathrm{step}+1}^V\) \textbf{and} \(w \notin \mathrm{path}\)}{
        \(\mathrm{path}.\texttt{push}(w)\)\;
        \(\mathrm{recurse}(w, \mathrm{step} + 1)\)\;
        \(\mathrm{path}.\texttt{pop}()\)\;
    }
}
\end{algorithm}

\subsection{Complejidad}

En el peor caso (sin podas), DFS subgraph match es:
\[
O\left(|\V_{L_0^V}| \cdot \prod_{i=0}^{k-1} \deg^{+}_{L_i^E, L_{i+1}^V}\right)
\]
donde \(|\V_{L_0^V}|\) es el número de vértices del tipo del primer
paso, y \(\deg^{+}_{L^E, L^V}\) es el grado promedio saliente con
\emph{esa} combinación etiqueta-arista, etiqueta-destino. Para
hipótesis cortas y etiquetas selectivas, ese producto es chico
($\sim 10^4$ para 3 saltos en datos típicos). Para hipótesis largas
o etiquetas no selectivas, explota.

\section{DFS temporal de GraphHunter}
\label{sec:dfs-temporal}

Sumando las restricciones temporales del capítulo~\ref{cap:grafos-temporales}:

\begin{algorithm}[H]
\caption{DFS subgraph match temporal de GraphHunter (versión simplificada)}
\KwIn{grafo \(G\), hipótesis \(H\) con ventana \(\Delta t\) y \(k_s\)-simplicidad}
\KwOut{Vec<HuntResult>}
\For{\(s \in \V\) con \(\mathrm{type}(s) = L_0^V\) \textbf{and} pasa \emph{prunes}}{
    \(\mathrm{path} \leftarrow [s]\)\;
    \(\mathrm{walk}(s, 0, \mathrm{None}, \mathrm{None})\)\;
}
\BlankLine
\Comment{\texttt{walk(v, step, t\_first, t\_prev)}}
\If{\(\mathrm{step} = k\)}{
    \texttt{emit}(\(\mathrm{path}\))\;
    \Return\;
}
\(t_{\mathrm{lo}} \leftarrow t_{\mathrm{prev}} + 1\) si \(t_{\mathrm{prev}}\) existe, sino \(-\infty\)\;
\(t_{\mathrm{hi}} \leftarrow t_{\mathrm{first}} + \Delta t\) si \(t_{\mathrm{first}}\) existe, sino \(+\infty\)\;
\For{\((v, w, \tau) \in \texttt{neighbors\_in\_window}(v, t_{\mathrm{lo}}, t_{\mathrm{hi}})\)}{
    \If{\(\mathrm{type}_E((v, w)) = L_{\mathrm{step}}^E\) \textbf{and} \(\mathrm{type}(w) = L_{\mathrm{step}+1}^V\) \textbf{and} \(\mathrm{count}(\mathrm{path}, w) < k_s\)}{
        \(\mathrm{path}.\texttt{push}(w)\)\;
        \(\mathrm{walk}(w, \mathrm{step}+1, t_{\mathrm{first}}\,\text{or}\,\tau, \tau)\)\;
        \(\mathrm{path}.\texttt{pop}()\)\;
    }
}
\end{algorithm}

\subsection{Las nuevas piezas}

\begin{itemize}
    \item \(t_{\mathrm{first}}\): timestamp de la \emph{primera}
    arista del walk. Se setea en el primer step y se usa para
    chequear ventana \(\Delta t\) de ahí en adelante.
    \item \(t_{\mathrm{prev}}\): timestamp del step anterior, para
    chequear monotonicidad estricta.
    \item \texttt{neighbors\_in\_window}: la operación clave del
    \emph{streaming backend} (capítulo~\ref{cap:streaming-mirror}).
    \texttt{streaming.neighbors\_arc(sid)} + \texttt{guard.read()} +
    \texttt{guard.as\_slice()} + dos llamadas a
    \texttt{partition\_point} devuelven el sub-slice contiguo de
    aristas salientes de \(v\) cuyo timestamp está en
    \([t_{\mathrm{lo}}, t_{\mathrm{hi}})\).
    \item \texttt{count(path, w) < k\_s}: \(k\)-simplicidad
    generalizada.
\end{itemize}

\subsection{Por qué iterativo y no recursivo en producción}

GraphHunter implementa el DFS de forma \emph{iterativa}
(\texttt{dfs\_match\_iterative}) usando una pila explícita en lugar
de la pila de recursión. Tres ventajas:

\begin{enumerate}
    \item \textbf{Sin stack overflow.} Hipótesis muy largas
    (k=10+) podrían reventar la pila de recursión (default 8~MB en
    Linux).
    \item \textbf{Cancelable.} El frame iterativo puede chequear
    una bandera de cancelación entre pasos sin requerir
    cooperación del compilador.
    \item \textbf{Cap más fácil.} Cuando el contador de resultados
    llega al cap, el loop iterativo simplemente sale; con
    recursión hay que propagar la señal hacia arriba.
\end{enumerate}

\section{De DFS naive a DFS bien-podado}

El algoritmo de la sección anterior es el \emph{baseline}. En la
práctica el motor agrega tres podas (\emph{prunes}) antes de
disparar el DFS principal:

\begin{enumerate}
    \item \textbf{NLF prune}: descarta \emph{starts} cuyo histograma
    de vecinos no satisface el primer paso
    (capítulo~\ref{cap:prune-nlf}).
    \item \textbf{k-hop reachability prune}: descarta starts desde
    los que no puede haber un path de longitud \(\leq
    \texttt{MAX\_DEPTH}\) al \(\mathrm{type}\) destino, dentro de
    la ventana temporal (capítulo~\ref{cap:prune-khop}).
    \item \textbf{Symmetry breaking}: para hipótesis palindrómicas
    (\(L_0^V = L_k^V\), \(L_0^E = L_{k-1}^E\), \dots), enumerar
    sólo la mitad y reflejar; ahorra factor \(2\times\).
\end{enumerate}

\noindent
Cada poda se aplica a los \emph{starts} antes del DFS principal,
o como early-exit dentro del walk. Las desarrollamos en los
próximos capítulos.

\section{Resumen}

\begin{itemize}
    \item DFS recursivo \emph{visit-all}: visita cada vértice
    alcanzable, marcado por \texttt{visited}. Lineal.
    \item DFS \emph{enumerate-paths}: enumera caminos sin marcar
    visited, usa backtracking, chequea \(k\)-simplicidad.
    \item DFS \emph{subgraph match}: agrega filtros de etiquetas
    (vértice y arista) en cada paso.
    \item DFS \emph{temporal}: agrega ventana \([t_{\mathrm{lo}},
    t_{\mathrm{hi}})\), monotonicidad, \texttt{neighbors\_in\_window}
    via búsqueda binaria.
    \item GraphHunter implementa la versión iterativa para
    cancelabilidad y stack-safety. La forma final está
    bien-podada (NLF, k-hop, symmetry breaking).
\end{itemize}

\begin{ejercicios}
\begin{enumerate}
    \item \starred{} Para el grafo de la figura~\ref{fig:dfs-ejemplo},
    enumerá todos los paths simples de longitud 3 desde \(A\). ¿Cuántos
    hay?

    \item Modificá el pseudocódigo de DFS \emph{enumerate-paths}
    para que devuelva la cantidad en vez de los paths. ¿Es más rápido?

    \item \starred{} Considerá DFS sobre un grafo
    \emph{Barabasi-Albert} de \(n = 10^4\), \(m = 10^5\), con
    grado max ($\Delta = 100$). Para una hipótesis de longitud 3
    sin filtros, ¿cuál es el peor caso de paths visitados?

    \item ¿Por qué la versión iterativa permite cancelación más
    fácil que la recursiva? Pensá la implementación de un
    \emph{cancellation token} en ambas formas.

    \item Diseñá un caso adversarial: un grafo y una hipótesis donde
    el DFS naive (sin podas) explote (\(> 10^9\) iteraciones).
    Identificá qué prune lo evitaría.
\end{enumerate}
\end{ejercicios}
